% Part: first-order-logic
% Chapter: introduction
% Section: sentences

\documentclass[../../../include/open-logic-section]{subfiles}

\begin{document}

\olfileid{fol}{int}{snt}

\olsection{\usetoken{P}{sentence}}

હવે આપણી પાસે !!{structure}s અને !!{formula}s માટે સંતોષ (``એમાં
સત્ય'')ની વ્યાખ્યાનો એક રૂપરેખાત્મક ખ્યાલ છે. પરંતુ તેમાં વધારાની એક
વસ્તુ---!!a{variable}-નિયુક્તિ---જોઈએ છે, જ્યારે આપણને તો
!!{sentence}sની વ્યાખ્યા જોઈતી હતી. નિયુક્તિઓને કેવી રીતે દૂર કરીશું
અને !!{sentence}s શું છે?

રૂપલક્ષી તર્કશાસ્ત્રના તમારા પ્રારંભિક અભ્યાસમાં !!{variable}s અને
પરિમાણકો વચ્ચેના સંબંધની ચર્ચા કદાચ યાદ હશે. પરિમાણક પછી હંમેશાં
!!a{variable} આવે છે અને પછી !!{sentence}ના જે ભાગને એ પરિમાણક લાગુ
પડે છે (તેનો ``વ્યાપ''), તેમાં એ !!{variable}ને પરિમાણકથી ``બદ્ધ''
સમજીએ છીએ. !!{formula}sમાં દરેક !!{variable} માટે તેને બાંધતો પરિમાણક
હોવો જરૂરી નહોતો; એવો પરિમાણક ન ધરાવતા !!{variable}s ``મુક્ત'' અથવા
``અબદ્ધ'' છે. આપણે મુક્ત !!{variable}s વિનાનાં બધાં !!{formula}sને
!!{sentence}s ગણીશું.

!!a{formula}~$!A$માં !!a{variable}નો કોઈ આવર્તન ક્યારે બદ્ધ છે, તેને
કયો પરિમાણક બાંધે છે અને તે ક્યારે મુક્ત છે તેનો સહજ ખ્યાલ ફરી સહેલો
છે. આ ચકાસવા તમે કદાચ કૌંસ ગણવાની કોઈ રીત શીખ્યા હશો. આપણે ચોક્કસ
વ્યાખ્યા આપવાનો આગ્રહ રાખવો પડે---અને !!{formula}sને અનુમાન વડે
વ્યાખ્યાયિત કર્યાં હોવાથી !!a{formula}~$!A$માં !!a{variable}~$x$ના
મુક્ત અને બદ્ધ આવર્તનો પણ અનુમાન વડે વ્યાખ્યાયિત કરી શકીએ. દાખલા
તરીકે, આપણી સરળ ભાષા માટે વ્યાખ્યા આવી દેખાઈ શકે:
\begin{enumerate}
  \item જો $!A$ આણ્વિક હોય, તો તેમાં $x$નાં બધાં આવર્તનો મુક્ત છે
  (એટલે કે,~$\Atom{\Obj P}{x}$માં $x$નું આવર્તન મુક્ત છે).
  \item જો $!A$નું સ્વરૂપ $\lnot !B$ હોય, તો~$\lnot !B$માં $x$નું
  આવર્તન ત્યારે અને માત્ર ત્યારે મુક્ત છે જ્યારે~$!B$માં $x$નું તેને
  અનુરૂપ આવર્તન મુક્ત હોય (એટલે કે,~$!B$માં ચલો જે મુક્ત આવર્તનોમાં
  આવે છે, એ જ અનુરૂપ આવર્તનો~$\lnot !B$માં મુક્ત છે).
  \item જો $!A$નું સ્વરૂપ $(!B \land !C)$ હોય, તો~$(!B \land !C)$માં
  $x$નું આવર્તન ત્યારે અને માત્ર ત્યારે મુક્ત છે જ્યારે~$!B$માં અથવા
  ~$!C$માં $x$નું તેને અનુરૂપ આવર્તન મુક્ત હોય.
  \item જો $!A$નું સ્વરૂપ $\lexists[x][!B]$ હોય, તો~$!A$માં $x$નું
  કોઈ આવર્તન મુક્ત નથી; જો તેનું સ્વરૂપ $\lexists[y][!B]$ હોય અને
  $y$,~$x$થી જુદું !!{variable} હોય, તો $\lexists[y][!B]$માં $x$નું
  આવર્તન ત્યારે અને માત્ર ત્યારે મુક્ત છે જ્યારે~$!B$માં $x$નું તેને
  અનુરૂપ આવર્તન મુક્ત હોય.
\end{enumerate}

ચલોનાં મુક્ત અને બદ્ધ આવર્તનોની ચોક્કસ વ્યાખ્યા મળી જાય પછી આપણે
સીધું કહી શકીએ: !!a{sentence} એટલે !!{variable}sનાં મુક્ત આવર્તનો
વિનાનું કોઈ પણ !!{formula}.

\end{document}
